\subsection{Homeomorphism groups of simply connected spaces}

Let X be a connected nonempty topological space, and let G be a discrete group acting continuously on the left on X; denote by e the identity element of G. Let M be a subset of X such that \(G \cdot M = X\) . We set

\[
\mathrm{S} = \{g \in \mathrm{G} \mid g \cdot \mathrm{M} \cap \mathrm{M} \neq \varnothing \}.
\]

We have \(e \in S\) and \(S = S^{-1}\) .

For every \(x \in X\) , denote by \(E_{x}\) the set of \(g \in G\) such that \(x \in g \cdot M\) . Let \(g, h \in E_{x}\) ; then \(g \cdot M \cap h \cdot M \neq \varnothing\) , so that \(g^{-1}h \in S\) . In particular, for every \(x \in M\) , we have \(e \in E_{x}\) whence \(E_{x} \subset S\) .

We make one of the following two hypotheses:

\begin{enumerate}
\def\labelenumi{(\roman{enumi})}
\item
  The set M is open;
\item
  The set M is closed and the covering \((g \cdot \mathrm{M})_{g \in \mathrm{G}}\) of X is locally finite.
\end{enumerate}

Lemma 2. --- For every point \(x \in X\) , the map \(\mu_x: E_x \times M \to X\) given by \((g, u) \mapsto g \cdot u\) is universally strict, and its image \(E_x \cdot M\) is a neighbourhood of \(x\) in \(X\) . In particular, \(S \cdot M\) is a neighbourhood of \(M\) .

Under assumption (i), the map \(\mu_{x}\) is open, since \(E_{x} \times M\) is open in \(G \times M\) . Its image is then open in X.

Suppose now that hypothesis (ii) is satisfied. The map \(\mu_x\) is proper (TG, I, p. 6, prop. 4, and p. 75, th. 1), hence universally strict (I, p. 20, corollary 11). Moreover, \((\mathrm{G} - \mathrm{E}_x) \cdot \mathrm{M}\) is a closed subset of X that does not contain \(x\) , so that \(\mathrm{E}_x \cdot \mathrm{M}\) is a neighbourhood of \(x\) .

The last assertion follows from the fact that \(E_{x} \subset S\) if \(x \in M\) .

Lemma 3. --- The group G is generated by S.

Let H be the subgroup of G generated by S. Let U = H·M ; observe that U is the union of the subsets of the form \(h \cdot (S \cdot M)\) , for \(h \in H\) . Since S · M is a neighbourhood of M (lemma 2), the set U is a neighbourhood of H · M = U ; it follows that U is open in X. Let \(g, g' \in G\) such that \(g \cdot U \cap g' \cdot U \neq \varnothing\) ; let \(h, h' \in H\) and \(x, x' \in M\) such that \(gh \cdot x = g'h' \cdot x'\) ; then \(h^{-1}g^{-1}g'h' \cdot x' = x\) , so that \(h^{-1}g^{-1}g'h' \in S\) ; in particular, \(g^{-1}g' \in H\) . When \(g\) runs through a system of representatives of the left cosets modulo \(H\) , the sets \(g \cdot U\) are thus pairwise disjoint; since \(G \cdot M = X\) , they cover \(X\) . Since \(X\) is connected, it follows that \((G : H) = 1\) , whence \(H = G\) .

Let T be the set of pairs \((s,t)\) of elements of G such that \(M\cap s\cdot M\cap st\cdot M\neq\varnothing\) ; if \((s,t)\in\mathrm{T}\) , then s, t, and st belong to S. Let F be the group \(\mathrm{F}(\mathrm{S},\mathbf{r})\) defined by the generating set S and by the set r of relations \(str^{-1}\) for \(r,s,t\in S\) such that \((s,t)\in\mathrm{T}\) and r=st; denote \(\varepsilon\) its identity element and \(\varphi\colon F\to G\) the canonical homomorphism (A, I, p. 86, prop. 9). If \(s\in S\) , we shall denote by \(x_{s}\) the element of F, image of s under the canonical map from S to F; for every \(s\in S\) , we have \(\varphi(x_{s})=s\) . The homomorphism \(\varphi\) is therefore surjective (lemma 3). For \((s,t)\in\mathrm{T}\) and r=st, we have \(x_{r}=x_{s}x_{t}\) ; one therefore has \(x_{e}=\varepsilon\) since \((e,e)\in\mathrm{T}\) ; for every \(s\in S\) , we also have \(x_{s^{-1}}=x_{s}^{-1}\) since \((s,s^{-1})\in\mathrm{T}\) .

Endow the set F with the discrete topology and denote by Z the topological space \(F \times M\) , equipped with the operation of F given by \((\gamma, (g, x)) \mapsto (\gamma g, x)\) , for \(\gamma\) and \(g \in F\) and \(u \in M\) . Let \((g_{1}, u_{1})\) and \((g_{2}, u_{2})\) be elements of Z; we shall say that \((g_{1}, u_{1})\) is congruent to \((g_{2}, u_{2})\) if there exists \(s \in S\) such that \(g_{2} = g_{1} \cdot x_{s}\) and \(s \cdot u_{2} = u_{1}\) .

Lemma 4. --- The relation “ \((g_{1}, u_{1})\) is congruent to \((g_{2}, u_{2})\) ” is an equivalence relation on Z, compatible with the operation of F.

This relation is reflexive, since we have \(x_e = \varepsilon\) . Let \((g_1, u_1)\) and \((g_2, u_2)\) be elements of Z such that \((g_1, u_1)\) is congruent to \((g_2, u_2)\) . Let \(s \in S\) such that \(g_2 = g_1 x_s\) and \(s \cdot u_2 = u_1\) ; since \(x_{s^{-1}} = x_s^{-1}\) , we have \(g_1 = g_2 x_{s^{-1}}\) and \(u_2 = s^{-1} \cdot u_1\) , therefore \((g_2, u_2)\) is congruent to \((g_1, u_1)\) ; consequently, this relation is symmetric. Finally, let us prove that it is transitive. Let \((g_1, u_1)\) , \((g_2, u_2)\) and \((g_3, u_3)\) be points of Z such that \((g_1, u_1)\) and \((g_2, u_2)\) are congruent, as well as \((g_2, u_2)\) and \((g_3, u_3)\) . Let \(s\) and \(t\) be in S such that \(g_2 = g_1 x_s\) and \(s \cdot u_2 = u_1\) on the one hand, and \(g_3 = g_2 x_t\) and \(t \cdot u_3 = u_2\) on the other hand. We have \(u_1 = s \cdot u_2 = st \cdot u_3\) , therefore \(u_1 \in M \cap s \cdot M \cap st \cdot M\) which shows that \((s, t)\) belongs to T and that \(st\) belongs to S. We then have \(g_3 = g_2 x_t = g_1 x_s x_t = g_1 x_{st}\) and \(u_1 = st \cdot u_3\) therefore, \((g_1, u_1)\) and \((g_3, u_3)\) are congruent. Thus, the relation "being congruent" is an equivalence relation on Z. It is compatible with the operation of F on Z.

Let Y be the quotient topological space of Z for the equivalence relation defined above. Denote \(\pi\colon\mathrm{Z}\to\mathrm{Y}\) the canonical map.

We also denote by \(q \colon \mathrm{Z} \to \mathrm{X}\) the map given by \((g, x) \mapsto \varphi(g) \cdot x\) ; it is surjective (lemma 3). By passing to the quotient, the map \(q\) induces a continuous and surjective map \(p \colon \mathrm{Y} \to \mathrm{X}\) ; by Lemma 4, the action of F on Z induces a continuous action of F on Y such that \(p(g \cdot y) = \varphi(g) \cdot p(y)\) for all \(g \in \mathrm{F}\) and every \(y \in \mathrm{Y}\) . In particular, the group \(\mathrm{N} = \mathrm{Ker}(\varphi)\) acts continuously on the X-space \((\mathrm{Y}, p)\) .

Lemma 5. --- If M is connected, then the space Y is connected.

Let \(g \in F\) and \(s \in S\) . Let u and v be elements of M such that \(v = s \cdot u\) ; one therefore has \(\pi(g, v) = \pi(gx_s, u)\) and the sets \(\pi(\{g\} \times M)\) and \(\pi(\{gx_s\} \times M)\) of Y have a point in common. Since they are connected, they are contained in the same connected component of Y. Since S is a symmetric subset of the group G and \(x_{s^{-1}} = x_s^{-1}\) for all \(s \in S\) , every element g of F is of the form \(x_{s_1} \ldots x_{s_n}\) , where \(n \in \mathbb{N}\) and \(s_1, \ldots, s_n\) are elements of S. By induction on n, the sets \(\pi(\{e\} \times M)\) and \(\pi(\{g\} \times M)\) are contained in the same connected component of Y, for every \(g \in F\) . It follows that Y is connected.

Lemma 6. --- For every \(x \in X\) , the group N acts faithfully and transitively on the fiber \(p^{-1}(x)\) .

Since \(p\) is surjective, the fiber \(p^{-1}(x)\) is nonempty. Let \(y, y' \in p^{-1}(x)\) ; let us prove that there exists a unique element \(n \in \mathrm{N}\) such that \(n \cdot y = y'\) . Let \(g, h \in \mathrm{F}\) and let \(u, v \in \mathrm{M}\) be such that \(y = \pi(g, u)\) and \(y' = \pi(h, v)\) . Set \(s = \varphi(g^{-1}h)\) . Since \(x = \varphi(g) \cdot u = \varphi(h) \cdot v\) , we have \(u = s \cdot v\) , whence \(s \in \mathrm{S}\) . It follows that \(\varphi(h) = \varphi(gx_s)\) , so that there exists \(n \in \mathrm{N}\) such that \(h = ngx_s\) . Then,

\[
y ^ {\prime} = \pi (h, v) = \pi (n g x _ {s}, v) = n \cdot \pi (g x _ {s}, v) = n \cdot \pi (g, s \cdot v) = n \cdot y.
\]

Let \(n'\) be an element of N such that \(n' \cdot y = y'\) . We have \(n' \cdot \pi(g, u) = n \cdot \pi(g, u)\) , whence \(\pi(n'g, u) = \pi(ng, u)\) . Consequently, there exists \(t \in S\) such that \(n'g = ngx_t\) ; this relation implies that \(\varphi(x_t) = e\) , whence \(t = e\) and \(n = n'\) .

Lemma 7. --- Equipped with the action of N, the X-space (Y,p) is a left principal covering. It is trivializable over M.

Let \(x \in \mathrm{X}\) ; we fix an element \(g \in \mathrm{E}_x\) as well as an element \(\overline{g} \in \mathrm{F}\) such that \(\varphi(\overline{g}) = g\) . We denote by \(\mu_x \colon \mathrm{E}_x \times \mathrm{M} \to \mathrm{X}\) the map given by \((h, u) \mapsto h \cdot u\) .

Let \(n \in \mathrm{N}\) , let \(h, k \in \mathrm{E}_x\) and let \(u, v \in \mathrm{M}\) be such that \(h \cdot u = k \cdot v\) ; set \(s = g^{-1}h\) , \(t = h^{-1}k\) and \(r = st = g^{-1}k\) . Since \(x\) belongs to \(g \cdot M \cap h \cdot M \cap k \cdot M\) , the pair \((s, t)\) belongs to T, hence \(x_s x_t = x_r\) . We therefore have

\[
\pi (n \overline {{g}} x _ {r}, v) = \pi (n \overline {{g}} x _ {s} x _ {t}, v) = \pi (n \overline {{g}} x _ {s}, t \cdot v) = \pi (n \overline {{g}} x _ {s}, u).
\]

Thus there exists a unique map \(\theta\colon\mathrm{N}\times(\mathrm{E}_{x}\cdot\mathrm{M})\to\mathrm{Y}\) such that \(\theta(n,h\cdot u)=\pi(n\overline{g}x_{g^{-1}h},u)\) for all \(n\in\mathrm{N}\) , all \(h\in\mathrm{E}_{x}\) and every \(u\in\mathrm{M}\) . We have

\[
p (\theta (n, h \cdot u)) = p (\pi (n \overline {{g}} x _ {s}, u)) = q (n \overline {{g}} x _ {s}, u) = \varphi (n \overline {{g}} x _ {s}) \cdot u = g s \cdot u = h \cdot u.
\]

Additionally, for \(n, n' \in \mathrm{N}\) and \(y \in \mathrm{E}_x \cdot \mathrm{M}\) , we have \(\theta(n'n, y) = n' \cdot \theta(n, y)\) . The map \(\theta'\colon \mathrm{N} \times (\mathrm{E}_x \times \mathrm{M}) \to \mathrm{Y}\) given by \(\theta'(n, (h, u)) = \theta(n, \mu_x(h, u))\) is continuous; since the map \(\mu_x\) is universally strict (I, p. 133, lemma 2), the map \(\theta\) is continuous. It is a bijection from \(\mathrm{N} \times (\mathrm{E}_x \cdot \mathrm{M})\) onto the subspace \(\mathrm{Y} \times_{\mathrm{X}} (\mathrm{E}_x \cdot \mathrm{M})\) of Y (Lemma 6).

Let \(z = (\overline{k}, v) \in \mathbf{Z}\) and \((h, u) \in \mathrm{E}_x \times \mathrm{M}\) such that \(q(\overline{k}, v) = \mu_x(h, u)\) . Set \(s = h^{-1}\varphi(\overline{k})\) ; since \(\varphi(\overline{k}) \cdot v = h \cdot u\) , we have \(s \in S\) . Let us then set \(\lambda'(z, (h, u)) = \overline{k} x_s^{-1} x_{h^{-1}g} \overline{g}^{-1}\) ; we have \(\varphi(\lambda'(z, (h, u))) = \varphi(\overline{k}) s^{-1} h^{-1} gg^{-1} = e\) , therefore \(\lambda'(z, (h, u)) \in \mathrm{N}\) . One thus defines a continuous map \(\lambda' \colon \mathbf{Z} \times_{\mathbf{X}} (\mathrm{E}_x \times \mathrm{M}) \to \mathrm{N}\) . Moreover, for every \(z\) and \((h, u)\) as above, one has

\[
\begin{array}{r l} & {\theta (\lambda^ {\prime} (z, (h, u)), h \cdot u) = \pi (\overline {{k}} x _ {s} ^ {- 1} x _ {h ^ {- 1} g} \overline {{g}} ^ {- 1} \overline {{g}} x _ {g ^ {- 1} h}, u)} \\ & {\qquad = \pi (\overline {{k}} x _ {s} ^ {- 1}, u) = \pi (\overline {{k}}, s ^ {- 1} \cdot u) = \pi (\overline {{k}}, v) = \pi (z),} \end{array}
\]

and \(\lambda'(z,(h,u))\) is the unique element \(n\) of N such that \(\theta(n,h\cdot u)=\pi(z)\) . There exists in particular a unique map

\[
\lambda \colon Y \times_ {X} (E _ {x} \cdot M) \to N
\]

such that \(\lambda(\pi(z), h \cdot u) = \lambda'(z, (h, u))\) for all \(z \in \mathbf{Z}\) and every \((h, u) \in \mathrm{E}_x \times \mathrm{M}\) such that \(q(z) = h \cdot u\) . It is continuous, because the map \(\mu_x\) is universally strict. It follows that the \(\mathrm{E}_x \cdot \mathrm{M}\) -space \(\mathrm{Y}_{\mathrm{E}_x \cdot \mathrm{M}}\) deduced from Y by base change, equipped with the operation of N, is a principal fiber bundle with group N, trivializable.

The lemma is thus proved.

PROPOSITION 10. --- Let X be a nonempty connected topological space, let G be a discrete group acting continuously on the left on X, and let M be a subset of X such that G · M = X. One makes one of the following two hypotheses:

\begin{enumerate}
\def\labelenumi{(\roman{enumi})}
\item
  The set M is open;
\item
  The set M is closed and the covering \((g \cdot \mathrm{M})_{g \in \mathrm{G}}\) of X is locally finite.
\end{enumerate}

Let S be the set of elements \(g \in G\) such that \(M \cap g \cdot M \neq \varnothing\) , let T be the set of pairs \((s, t) \in S \times S\) such that \(M \cap s \cdot M \cap st \cdot M \neq \varnothing\) . Let F be the group \(F(S, r)\) defined by the generating set S and by the set r of relations \(str^{-1}\) for \(r, s, t \in S\) such that \((s, t) \in T\) and r = st; for \(s \in S\) , denote by \(x_s\) the image of s under the canonical map from S to F. There exists a unique group homomorphism \(\varphi: F \to G\) such that \(\varphi(x_s) = s\) for all \(s \in S\) . It is surjective; it is an isomorphism if the space X is simply connected, or, more generally, if every covering of X that is trivializable over M is trivializable.

The homomorphism \(\varphi\) is surjective (I, p. 133, lemma 3). With the notation just introduced, the covering Y of X is trivializable, since it is trivializable over M. By lemma 5, Y is connected. Consequently, \(p\colon\mathrm{Y}\to\mathrm{X}\) is a homeomorphism and the group N is reduced to the identity element. Consequently, the homomorphism \(\varphi\colon\mathrm{F}\to\mathrm{G}\) is a group isomorphism.

PROPOSITION 11. --- Let X be a simply connected topological space and G a discrete group acting continuously on the right on X. If the subgroup of G generated by the union of the stabilizers of the points of X is equal to G, then the space X/G is simply connected.

Let (E, p) be a covering of X/G. We need to show that, for every point x of E, there exists a continuous section s of p such that \(s \circ p(x) = x\) (I, p. 70, cor. 2 of prop. 1). Let us denote by \(q: X \to X/G\) the canonical map, and choose a point y of X such that \(q(y) = p(x)\) . Since the space X is simply connected, there exists a continuous map \(f: X \to E\) such that \(f(y) = x\) and \(p \circ f = q\) (I, p. 125, prop. 3). Let z be a point of X and g an element of the subgroup of G fixing z. The maps \(t \mapsto f(t)\) and \(t \mapsto f(t \cdot g)\) from X into E are continuous lifts to E of the map q that coincide for t = z. Since the space X is connected, they are equal (I, p. 34, cor. 1 of prop. 11). Since the union of the stabilizer subgroups of the points of X generates G, we have \(f(t \cdot g) = f(t)\) for all \(t \in X\) and every \(g \in G\) . By passing to the quotient, we deduce from f a continuous map s: X/G → E which is a continuous section of p such that \(s(p(x)) = x\) .

